\scanpage{108}
c'est aussi une section rationnelle $\varphi$ de $L'$, et en fait une \deleted{section régulière} \inserted{\unsure{pseudo-section}} \deleted{[car \unsure{on se ramène à} $X'$ intègre] \unsure{d'origine de} $\eta$}; soit $D'$ son diviseur.
\inserted{qui d'ailleurs n'est déterminée (\unsure{$L'$ défini|$L'$ fixé}) que mod multiplication par un $f\in k(\eta)^*$.}

Alors $D'$ est un diviseur de \uline{Cartier} \deleted{et} sur $X'$, \deleted{et on} je dis qu'il est \uline{transversal aux fibres}. \inserted{(quitte à multiplier $\varphi$ par $f\in k(\eta)^*$)} En effet, comme $X'_{s'}$ est intègre, il suffit de voir que
\[
 \operatorname{supp}D'\not\supset X'_{s'}.
\]
Or sinon \deleted{un} $D'$ contient \deleted{cette} \unsure{composante} \deleted{de} $X'_{s'}$, \deleted{soit} comme $X'_{s'}$ est intègre, $D'-nX'_{s'}$ \deleted{sera} n'aura plus la composante $X'_{s'}$. Il suffit donc de remplacer $\varphi$ par $\pi^{-n}\varphi$, où $\pi$ est l'uniformisante de $V$.

\margin{\uline{Variante} On normalise $\varphi$ par la condition que $\varphi$ coïncide avec la section canonique de $L'$ le long de la section unité.}

Donc\ldots\ $L'_{s'}\simeq L(D'_{s'})$, d'autre part
\[
 \operatorname{supp}D'_{s'}\subset\bigl(\overline{\operatorname{supp}D'_{\eta'}}\bigr)_{s'}\deleted{\unsure{Z'_{s'}}}
\]
[\uline{vrai pour $\forall$ diviseur de Cartier transversal aux fibres}. \unsure{Est ramené au cas} $\geq0$, \unsure{où c'est platitude}\ldots].
\margin{lemme facile}
\margin{107}
